\documentclass[11pt,a4paper]{article}

%% ─── Gói lệnh ────────────────────────────────────────────────────────────────
\usepackage[a4paper, margin=2.5cm]{geometry}
\usepackage[utf8]{inputenc}
\usepackage[T5]{fontenc}
\usepackage{lmodern}
\usepackage[vietnamese,english]{babel}
\usepackage{microtype}
\usepackage{amsmath,amssymb}
\usepackage{bm}
\usepackage{booktabs}
\usepackage{array}
\usepackage{multirow}
\usepackage{graphicx}
\usepackage{listings}
\usepackage{xcolor}
\usepackage{hyperref}
\usepackage{caption}
\usepackage{subcaption}
\usepackage{enumitem}
\usepackage{fancyhdr}
\usepackage{titlesec}
\usepackage{parskip}
\usepackage{placeins}

%% ─── Hyperref ────────────────────────────────────────────────────────────────
\hypersetup{
    colorlinks=true,
    linkcolor=blue!70!black,
    citecolor=green!50!black,
    urlcolor=blue!60!black,
    pdftitle={128T128R NZP-CSI-RS: Tu Nguyen Mau Don Khoi den Danh Gia Quet SNR},
}

%% ─── Đầu trang / Chân trang ──────────────────────────────────────────────────
\pagestyle{fancy}
\fancyhf{}
\lhead{\small Báo cáo Mô phỏng 128T128R NZP-CSI-RS}
\rhead{\small ThangTQ23 -- VSI}
\cfoot{\thepage}

%% ─── Định dạng code MATLAB ───────────────────────────────────────────────────
\lstdefinestyle{matlab}{
    language=Matlab,
    basicstyle=\ttfamily\small,
    keywordstyle=\color{blue!80!black}\bfseries,
    commentstyle=\color{green!50!black}\itshape,
    stringstyle=\color{orange!80!black},
    numbers=left,
    numberstyle=\tiny\color{gray},
    numbersep=8pt,
    frame=single,
    framesep=4pt,
    breaklines=true,
    tabsize=4,
    showstringspaces=false,
    captionpos=b,
}

%% ─── Tiêu đề ─────────────────────────────────────────────────────────────────
\title{
    \vspace{-1.5cm}
    \Large\textbf{Mô phỏng NZP-CSI-RS hệ thống 128T128R:}\\[4pt]
    \large Từ Nguyên mẫu Đơn khối (v1.0) đến\\
    Đánh giá Quét SNR Mô-đun hóa (v2.0)\\[8pt]
    \normalsize Báo cáo Kỹ thuật
}
\author{ThangTQ23 \\ VSI}
\date{Tháng 3 năm 2026}

%% ─── Tài liệu ────────────────────────────────────────────────────────────────
\begin{document}

\maketitle
\thispagestyle{fancy}

\begin{abstract}
Tài liệu này mô tả quá trình xây dựng công cụ mô phỏng phản hồi CSI cho hệ thống
Massive MIMO 128T128R, dùng NZP-CSI-RS Row~18 với CDM8. Điểm xuất phát là v1.0
-- một script MATLAB chạy một lần, một điểm SNR, bốn phương pháp CSI gộp chung
vào một file. Qua hai giai đoạn phát triển, v2.0 được tái cấu trúc thành pipeline
mô-đun, bổ sung Monte Carlo và quét SNR, đồng thời nâng cấp cách tính MCS sang
chuẩn L2SM/MIESM. Kết quả cho thấy Mode~B (Phương pháp~D) đạt 33--70\,\% so với
cận trên SVD tùy SNR, và luôn nhỉnh hơn Mode~A (Phương pháp~C) khoảng 5--6\,\%.
\end{abstract}

\tableofcontents
\newpage

%% ═══════════════════════════════════════════════════════════════════════════
\section{Bối cảnh và Động lực}
%% ═══════════════════════════════════════════════════════════════════════════

\subsection{Tổng quan Hệ thống}

Hệ thống mô phỏng nhắm đến cấu hình gNB \textbf{128T128R} ở băng Sub-6~GHz.
Mảng anten gNB là panel $4\text{V}\times16\text{H}\times2\text{pol} = 128$ cổng,
UE dùng mảng $1\text{V}\times2\text{H}\times2\text{pol} = 4$ cổng.
Câu hỏi cốt lõi là: codebook Rel-19 Type~I Single-Panel 128 cổng (Mode~A và Mode~B)
đạt được bao nhiêu phần trăm so với cận trên SVD lý tưởng, ở các mức SNR khác nhau?

\subsection{Vì sao cần v2.0?}

V1.0 chạy được pipeline đầu cuối nhưng chỉ ở một điểm SNR và một realization kênh.
Như vậy kết quả còn mang nặng tính ngẫu nhiên, chưa đại diện thống kê. V2.0 giải
quyết các điểm đó:

\begin{itemize}[noitemsep]
    \item Thêm Monte Carlo: lấy trung bình RI/MCS/TP qua nhiều realization kênh
          để kết quả có ý nghĩa thống kê.
    \item Quét 3 điểm SNR ($-5$, $+10$, $+25$\,dB) tương ứng ô xa/trung bình/gần.
    \item Thêm cột \% so với SVD để dễ so sánh Mode~A và Mode~B.
    \item Sửa cách tính MCS: dùng L2SM/MIESM thay vì dùng thẳng chỉ số CQI.
    \item Tách code thành các file hàm riêng, dễ test và tái sử dụng.
\end{itemize}

%% ═══════════════════════════════════════════════════════════════════════════
\section{Tham số Mô phỏng}
%% ═══════════════════════════════════════════════════════════════════════════

\subsection{Cấu hình Sóng mang}

Có hai chế độ mô phỏng tùy mục đích: chế độ nhanh (10\,MHz) dùng khi cần chạy
thử nhanh nhiều lần, chế độ chuẩn (100\,MHz) dùng khi cần kết quả sát thực tế hơn.

\begin{table}[h!]
\centering
\caption{Tham số Sóng mang 5G NR}
\label{tab:carrier}
\begin{tabular}{llll}
\toprule
\textbf{Tham số} & \textbf{Mô phỏng nhanh} & \textbf{Chuẩn} & \textbf{Ghi chú} \\
\midrule
Subcarrier Spacing & 30\,kHz  & 30\,kHz   & FR1 NR \\
NSizeGrid          & 52 PRB   & 273 PRB   & 10\,MHz / 100\,MHz \\
CyclicPrefix       & Normal   & Normal    & \\
\bottomrule
\end{tabular}
\end{table}

\subsection{Cấu hình Anten}

\begin{table}[h!]
\centering
\caption{Cấu hình Mảng Anten}
\label{tab:antennas}
\begin{tabular}{lll}
\toprule
\textbf{Nút} & \textbf{Kích thước mảng} & \textbf{Tổng số cổng} \\
\midrule
gNB (TX) & $4\text{V}\times16\text{H}\times2\text{pol}$ & 128 \\
UE (RX)  & $1\text{V}\times2\text{H}\times2\text{pol}$  & 4 \\
\bottomrule
\end{tabular}
\end{table}

\subsection{Cấu hình NZP-CSI-RS (Row~18, CDM8)}

Để phủ 128 cổng TX, bốn tài nguyên NZP-CSI-RS được sử dụng, mỗi tài nguyên mang
32 cổng qua Row~18 (CDM8 = FD-CDM2~$\times$~TD-CDM4). Các tài nguyên được
trải đều trên hai slot liên tiếp để tránh va chạm symbol:

\begin{table}[h!]
\centering
\caption{Bố cục CSI-RS 4 tài nguyên trên 2 slot}
\label{tab:csirs}
\begin{tabular}{ccccc}
\toprule
\textbf{Tài nguyên} & \textbf{Slot} & \textbf{Symbol bắt đầu ($\ell_0$)} &
\textbf{Các symbol} & \textbf{Cổng} \\
\midrule
R0 & 0 & 2 & 2, 3, 4, 5   & 0--31  \\
R1 & 0 & 8 & 8, 9, 10, 11 & 32--63 \\
R2 & 1 & 2 & 2, 3, 4, 5   & 64--95 \\
R3 & 1 & 8 & 8, 9, 10, 11 & 96--127 \\
\bottomrule
\end{tabular}
\end{table}

Trải phổ CDM8 sử dụng FD-CDM độ dài~2 và TD-CDM độ dài~4, tạo ra các chuỗi pilot
trực giao trên $2\times4 = 8$ nhóm CDM, mỗi nhóm mang 4 cổng, tổng cộng 32 cổng
mỗi tài nguyên. Vị trí sóng mang con theo TS~38.211 Bảng~7.4.1.5.3-1, Row~18,
mật độ~0.5: $k' \in \{0, 2, 4, 6\}$.

\subsection{Mô hình Kênh}

Tương tự sóng mang, tham số kênh cũng có hai chế độ. Chế độ nhanh dùng delay spread
thấp và Doppler nhỏ để kênh ổn định, dễ ước lượng, giúp kiểm tra pipeline dễ hơn.
Chế độ chuẩn dùng tham số sát với môi trường thực tế Urban Macro 4{,}9\,GHz.

\begin{table}[h!]
\centering
\caption{Tham số Kênh CDL}
\label{tab:channel}
\begin{tabular}{llll}
\toprule
\textbf{Tham số} & \textbf{Mô phỏng nhanh} & \textbf{Chuẩn} & \textbf{Ghi chú} \\
\midrule
Delay Profile        & CDL-B              & CDL-B              & \\
Delay Spread         & 100\,ns            & 450\,ns            & Urban Micro / Macro \\
Doppler tối đa       & 5\,Hz              & 136\,Hz            & 1\,km/h / 30\,km/h @ 4{,}9\,GHz \\
Tần số sóng mang     & \multicolumn{2}{c}{4{,}9\,GHz}        & \\
Khoảng cách phần tử  & \multicolumn{2}{c}{0{,}5\,$\lambda$}  & TX và RX \\
\bottomrule
\end{tabular}
\end{table}

\noindent\textbf{Lưu ý:} CDL-B ở delay spread thấp / Doppler nhỏ cho kênh rank-4
khá đều, phù hợp để kiểm tra codebook đa lớp. Khi tăng delay spread lên 450\,ns
và Doppler lên 136\,Hz, chất lượng ước lượng kênh giảm do trung bình $\mathbf{H}_{wb}$
bị trải rộng hơn trong miền tần số và thời gian -- điều này làm RI giảm, nhưng đây
là hiệu ứng của kênh, không phải do codebook bị giới hạn.

\subsection{Tham số Quét Mô phỏng}

\begin{table}[h!]
\centering
\caption{Cấu hình Quét SNR Monte Carlo}
\label{tab:sweep}
\begin{tabular}{ll}
\toprule
\textbf{Tham số} & \textbf{Giá trị} \\
\midrule
Các điểm SNR & $\{-5, +10, +25\}$\,dB (xa / trung bình / gần) \\
Số realization mỗi SNR & 20 \\
Tổng số lần chạy & 60 \\
\bottomrule
\end{tabular}
\end{table}

%% ═══════════════════════════════════════════════════════════════════════════
\section{Kiến trúc Phiên bản 1.0}
%% ═══════════════════════════════════════════════════════════════════════════

Phiên bản 1.0 (\texttt{CSIRS\_128T128R\_v1\_0.m}) là một script MATLAB đơn khối
được tổ chức thành mười hai mục:

\begin{enumerate}[noitemsep]
    \item Cấu hình hệ thống (sóng mang, anten, tham số)
    \item Cấu hình 4 tài nguyên NZP-CSI-RS (Row~18, CDM8)
    \item Tạo symbol CSI-RS và ánh xạ lên lưới tài nguyên 2 slot
    \item Trực quan hóa lưới tài nguyên CSI-RS
    \item Thiết lập kênh CDL
    \item Tạo dạng sóng TX và điều chế OFDM
    \item Truyền qua kênh, AWGN, ước lượng định thời, giải điều chế
    \item Ước lượng kênh (mỗi tài nguyên dùng \texttt{nrChannelEstimate})
    \item Tính NMSE (so với kiến thức kênh hoàn hảo)
    \item Phản hồi CSI: Phương pháp A, B, C, D
    \item Trực quan hóa đồ hình Beamforming
    \item Bảng tóm tắt kết quả
\end{enumerate}

\subsection{Định nghĩa các Phương pháp (v1.0)}

\begin{description}[style=nextline]
    \item[Phương pháp A -- Rel-15/16 Type~I Single Panel mỗi 32 cổng]
        Mỗi trong bốn tài nguyên 32 cổng được xử lý độc lập bằng codebook
        chuẩn Rel-15/16 (\texttt{typeI-SinglePanel}). PMI từ từng mảng con được
        kết hợp để tạo precoder tổng hợp. Phương pháp này hoạt động trực tiếp
        với thư viện MATLAB 5G chuẩn mà không cần sửa đổi.
    \item[Phương pháp B -- Cận trên SVD đầy đủ 128 cổng]
        Kênh wideband $\mathbf{H}_{wb} \in \mathbb{C}^{N_r \times 128}$ được phân
        tích bằng SVD. Precoder tối ưu $\mathbf{W}_B$ sử dụng $r_B$ vector kỳ dị
        phải hàng đầu. Đây là tham chiếu lý tưởng (không có tổn hao lượng hóa codebook).
    \item[Phương pháp C -- Rel-19 Type~I Single-Panel Mode~A]
        Sử dụng codebook Rel-19 đầy đủ 128 cổng ở chế độ beam DFT oversample
        (\texttt{typeI-SinglePanel-r19}, \texttt{CodebookMode=1}).
    \item[Phương pháp D -- Rel-19 Type~I Single-Panel Mode~B]
        Sử dụng codebook Rel-19 ở chế độ cơ sở trực giao
        (\texttt{typeI-SinglePanel-r19}, \texttt{CodebookMode=2}), cung cấp độ
        phân giải không gian tốt hơn với chi phí overhead phản hồi cao hơn.
\end{description}

\subsection{Hạn chế của v1.0}

\begin{itemize}[noitemsep]
    \item Chỉ một điểm SNR (\texttt{SNRdB}~= 20\,dB) --- không có quét bán kính ô.
    \item Chỉ một realization kênh --- kết quả không có ý nghĩa thống kê.
    \item CQI tính từ bảng ngưỡng cố định trên SINR ZF trung bình;
          MCS lấy trực tiếp từ chỉ số CQI (ánh xạ không đúng chuẩn).
    \item Phương pháp~A (Rel-15/16) được tích hợp nhưng ít giá trị cho đánh giá Rel-19.
    \item Toàn bộ 12 mục trong một file --- khó tái sử dụng, kiểm thử hay mở rộng.
\end{itemize}

%% ═══════════════════════════════════════════════════════════════════════════
\section{Phiên bản 2.0: Kiến trúc Mô-đun}
%% ═══════════════════════════════════════════════════════════════════════════

\subsection{Nguyên tắc Thiết kế}

V2.0 phân tách script đơn khối thành một file chính và bảy file hàm,
mỗi file có một trách nhiệm được xác định rõ ràng. Sự phân tách này cho phép:

\begin{itemize}[noitemsep]
    \item Kiểm thử đơn vị độc lập cho từng giai đoạn pipeline.
    \item Tái sử dụng lưới TX (chỉ xây dựng một lần, ngoài tất cả các vòng lặp)
          cho toàn bộ các tổ hợp SNR/realization.
    \item Luồng dữ liệu rõ ràng: mỗi hàm chỉ nhận đầu vào mà nó cần và trả về
          đúng những gì giai đoạn tiếp theo yêu cầu.
\end{itemize}

\subsection{Cấu trúc File}

\begin{table}[h!]
\centering
\caption{Cấu trúc File V2.0}
\label{tab:files}
\begin{tabular}{lp{9cm}}
\toprule
\textbf{File} & \textbf{Trách nhiệm} \\
\midrule
\texttt{CSIRS\_128T128R\_v2\_0.m}    & File chính: cấu hình hệ thống, vòng lặp quét SNR, Monte Carlo, tổng hợp kết quả \\
\texttt{csirs\_buildGrid.m}          & Xây dựng đối tượng CSI-RS và lưới TX 2 slot (chạy một lần) \\
\texttt{csirs\_channelSetup.m}       & Tạo và cấu hình đối tượng \texttt{nrCDLChannel} \\
\texttt{csirs\_txRx.m}               & Điều chế OFDM, truyền qua kênh, AWGN, giải điều chế \\
\texttt{csirs\_channelEstimate.m}    & \texttt{nrChannelEstimate} từng tài nguyên, lắp ráp $\mathbf{H}_{est}^{full}$ \\
\texttt{csirs\_feedback.m}           & Phản hồi CSI: RI/PMI/CQI cho Phương pháp B, C, D \\
\texttt{csirs\_computeMetrics.m}     & RI, MCS (L2SM), Thông lượng từ struct phản hồi \\
\texttt{csirs\_printSweepTable.m}    & In bảng ASCII với \% so với tham chiếu SVD \\
\bottomrule
\end{tabular}
\end{table}

\FloatBarrier
\subsection{Cấu trúc Vòng lặp Chính}

Script chính v2.0 thực hiện vòng lặp hai cấp:

\begin{lstlisting}[style=matlab, caption={V2.0 two-level simulation loop (simplified)}]
% One-time setup (outside all loops)
[csirs, txGrid_2slots, ...] = csirs_buildGrid(carrier, nTxAntennas);
channel = csirs_channelSetup(carrier, gnbArraySize, ueArraySize, channelCfg);

for iSNR = 1 : nSNR                      % Outer: SNR points
    for iReal = 1 : nRealiz              % Inner: Monte Carlo
        release(channel);                % Draw new independent channel realization

        % RX chain
        [rxGrid_s0, rxGrid_s1, ...] = csirs_txRx(..., SNRdB, ...);
        [H_est_full, nVar_all]       = csirs_channelEstimate(...);

        % CSI chain
        fb = csirs_feedback(carrier, csirs, H_est_full, nVar_all, slotAssign);
        m  = csirs_computeMetrics(fb, carrier);

        acc_ri(iReal,:)  = m.ri;
        acc_mcs(iReal,:) = m.mcs;
        acc_tp(iReal,:)  = m.tp;
    end
    results(iSNR).ri  = mean(acc_ri,  1);
    results(iSNR).mcs = mean(acc_mcs, 1);
    results(iSNR).tp  = mean(acc_tp,  1);
end
csirs_printSweepTable(results, snrList);
\end{lstlisting}

\paragraph{Tại sao \texttt{release(channel)} thực hiện được Monte Carlo.}
Đối tượng \texttt{nrCDLChannel} lưu trữ trạng thái nội bộ (seed ngẫu nhiên, bộ nhớ
bộ lọc, độ lệch định thời). Gọi \texttt{release()} sẽ đặt lại toàn bộ trạng thái
nội bộ và tạo một realization mới từ mô hình kênh ngẫu nhiên nền. Mỗi lần gọi
\texttt{release()} tiếp theo truyền dạng sóng là tương đương thống kê với việc
xây dựng một đối tượng kênh hoàn toàn mới, độc lập. Cách này vừa hiệu quả hơn
vừa tạo ra cùng phân phối như khi tạo $N_{\text{real}}$ kênh độc lập riêng biệt.

\subsection{Ước lượng Kênh và Tổng hợp Wideband}

Ước lượng kênh sử dụng \texttt{nrChannelEstimate} riêng biệt cho mỗi trong bốn
tài nguyên CSI-RS 32 cổng, với \texttt{CDMLengths}~= $[2,\,4]$ khớp với trải phổ
FD-CDM2 và TD-CDM4. Các ước lượng theo tài nguyên được lắp ráp thành ma trận
wideband đầy đủ:

\begin{equation}
\mathbf{H}_{wb} \in \mathbb{C}^{N_r \times N_t}
= \bigoplus_{r=0}^{3} \left\langle \hat{\mathbf{H}}^{(r)} \right\rangle_{k,\,\ell}
\end{equation}

trong đó $\langle\cdot\rangle_{k,\,\ell}$ ký hiệu trung bình theo chỉ số sóng
mang con và symbol OFDM trong slot của tài nguyên.

%% ═══════════════════════════════════════════════════════════════════════════
\FloatBarrier
\section{Pipeline Phản hồi CSI}
%% ═══════════════════════════════════════════════════════════════════════════

\subsection{Phương pháp B: Cận trên SVD}

Kênh wideband được phân tích:
$\mathbf{H}_{wb} = \mathbf{U}\mathbf{S}\mathbf{V}^H$.
RI được chọn theo:
\begin{equation}
r_B = \min\Bigl(\bigl|\{i : \sigma_i > 0{,}1\,\sigma_1\}\bigr|,\; N_r\Bigr)
\end{equation}
trong đó $\sigma_i$ là các giá trị kỳ dị. Precoder $\mathbf{W}_B = \mathbf{V}[:,\,1{:}r_B]$.

Tách kênh ZF được áp dụng cho kênh hiệu dụng
$\mathbf{H}_{eff} = \mathbf{H}_{wb}\mathbf{W}_B$:

\begin{equation}
\text{SINR}_\ell^{(B)}
= \frac{|\mathbf{w}_\ell^H\,\mathbf{h}_{eff,\ell}|^2}
       {\sum_{j \neq \ell}|\mathbf{w}_\ell^H\,\mathbf{h}_{eff,j}|^2
        + \sigma_n^2\|\mathbf{w}_\ell\|^2}
\end{equation}

trong đó $\mathbf{w}_\ell^H$ là hàng thứ $\ell$ của $(\mathbf{H}_{eff})^+$.

Dung lượng được tính:
\begin{equation}
C_B = \sum_{\ell=1}^{r_B} \log_2\!\left(1 + \frac{\sigma_\ell^2}{r_B\,\sigma_n^2}\right)
\;\text{[bits/s/Hz]}
\end{equation}

\subsection{Phương pháp C và D: Codebook Rel-19}

Cả hai phương pháp sử dụng \texttt{nr5g.internal.nrPMIReport} với
\texttt{CodebookType = 'typeI-SinglePanel-r19'} và
\texttt{PanelDimensions = [1, 16, 4]} (128 cổng: $N_1=16$, $N_2=4$).
RI được chọn qua tối đa hóa dung lượng:

\begin{equation}
r_{C/D}^* = \arg\max_{r \in \{1,\ldots,r_{\max}\}}\;
\log_2 \det\!\left(\mathbf{I} + \frac{1}{\sigma_n^2}\,
\mathbf{H}_{wb}\mathbf{W}(r)\mathbf{W}(r)^H\mathbf{H}_{wb}^H\right)
\end{equation}

trong đó $\mathbf{W}(r)$ là codeword wideband tốt nhất cho rank $r$ từ
\texttt{nrPMIReport}. Mode~A dùng \texttt{CodebookMode=1};
Mode~B dùng \texttt{CodebookMode=2}.

\subsection{CQI qua L2SM/MIESM (nâng cấp v2.0)}

Trong v1.0, CQI được tính từ bảng ngưỡng SINR cố định --- tương đương tra cứu
SINR-sang-CQI trực tiếp, không xét phân bố SINR theo RE hay ràng buộc BLER.

Trong v2.0, CQI thu được qua \texttt{nr5g.internal.nrCQISelect}, bên trong
thực hiện \textbf{Ánh xạ Link-to-System (L2SM)} bằng
\textbf{Ánh xạ SINR Hiệu dụng Thông tin Tương hỗ (MIESM)}:

\begin{equation}
\text{MI}_k = \log_2(1 + \text{SINR}_k), \quad
\overline{\text{MI}} = \frac{1}{N_{RE}}\sum_{k=1}^{N_{RE}}\text{MI}_k, \quad
\text{SINR}_{eff} = 2^{\overline{\text{MI}}} - 1
\end{equation}

SINR hiệu dụng được ánh xạ lên đường cong BLER (từ trừu tượng hóa link-level);
chỉ số MCS cao nhất thỏa mãn BLER~$\leq 10\%$ được chọn và chuyển sang CQI theo
TS~38.214~S5.2.2.1.

\subsection{MCS qua \texttt{computeMCS} (nâng cấp v2.0)}

Trong v1.0, MCS được xấp xỉ bằng cách dùng trực tiếp chỉ số CQI làm chỉ số MCS ---
một cách tắt phổ biến nhưng không đúng chuẩn. Trong v2.0,
\texttt{nr5g.internal.computeMCS} được gọi với ma trận SINR theo RE làm đầu vào,
áp dụng L2SM bên trong:

\begin{lstlisting}[style=matlab, caption={MCS computation via L2SM in csirs\_computeMetrics.m}]
l2sm   = nr5g.internal.L2SM.initialize(carrier);
pdsch  = makePDSCH(carrier, ri);              % nrPDSCHConfig: PRBSet + NumLayers
% Approach B: replicate per-layer wideband SINR to [nRE x ri]
sinr_B_mat = repmat(fb.sinr_B(:).', nPRB*12, 1);
mcs_B = nr5g.internal.computeMCS(l2sm, carrier, pdsch_B, sinr_B_mat, 'qam64');
% Approaches C/D: sinrPerRE is already [nRE x nLayers] from nrCQISelect
mcs_C = nr5g.internal.computeMCS(l2sm, carrier, pdsch_C, fb.sinrPerRE_C, 'qam64');
\end{lstlisting}

Thông lượng sau đó tính theo TS~38.214 Bảng~5.1.3.1-1:

\begin{equation}
\text{TP} = \frac{r \cdot Q_m \cdot R \cdot N_{RE/\text{slot}}}{T_{\text{slot}}}
\;\text{[bps]}
\end{equation}

trong đó $r$ là RI, $Q_m$ bậc điều chế, $R$ tốc độ mã,
$N_{RE/\text{slot}} = N_{PRB}\times12\times12$ (12 symbol PDSCH không có DMRS),
và $T_{\text{slot}} = 0{,}5\,\text{ms}$ tại SCS~30\,kHz.

%% ═══════════════════════════════════════════════════════════════════════════
\FloatBarrier
\section{Mở rộng Thư viện: \texttt{nrPMIReport} cho Rel-19}
%% ═══════════════════════════════════════════════════════════════════════════

Hàm \texttt{nr5g.internal.nrPMIReport} trong bản sao thư viện cục bộ đã hỗ trợ
Rel-15/16 nhưng cần được mở rộng để xử lý codebook \texttt{'typeI-SinglePanel-r19'}
và số cổng 128. Tổng cộng \textbf{5 thay đổi} trong \texttt{nrPMIReport.m} và
\textbf{2 hàm mới} đã được thêm vào.

\subsection{Lý thuyết Codebook Rel-19}

Cả Mode~A và Mode~B đều xây dựng bộ tiền mã hoá dạng dual-polarisation DFT.
Cấu trúc chung cho hạng $v$:

\begin{equation}
W^{(v)} = \frac{1}{\sqrt{v\,P}}
\begin{bmatrix}
  v_{l_1,m_1} & \cdots & v_{l_v,m_v} \\
  \varphi(c_1)\,v_{l_1,m_1} & \cdots & \varphi(c_v)\,v_{l_v,m_v}
\end{bmatrix}
\label{eq:precoder_general}
\end{equation}

với $P=128$, $\varphi(c) = e^{j\pi c/2}$, và vector beam DFT 2D:

\begin{equation}
v_{l,m} = \mathbf{u}_h(l) \otimes \mathbf{u}_v(m), \quad
[\mathbf{u}_h(l)]_k = e^{j2\pi lk/(O_1 N_1)}, \quad
[\mathbf{u}_v(m)]_k = e^{j2\pi mk/(O_2 N_2)}
\end{equation}

với $(N_1,N_2)=(16,4)$, $(O_1,O_2)=(4,4)$ cho cấu hình 128 cổng.

\medskip\noindent\textbf{Điểm khác nhau giữa hai mode} là cách xác định chỉ số
beam $(l_k, m_k)$ của mỗi tầng:

\begin{table}[h!]
\centering
\caption{So sánh cấu trúc chỉ số Mode A vs Mode B}
\label{tab:modeAB_diff}
\begin{tabular}{p{3.5cm}p{5cm}p{5cm}}
\toprule
 & \textbf{Mode A (\texttt{CodebookMode=1})} & \textbf{Mode B (\texttt{CodebookMode=2})} \\
\midrule
Chỉ số beam & $l \in [0,N_1O_1{-}1]$, $m \in [0,N_2O_2{-}1]$ (chung) & $(q_1,q_2)$ offset + $i_{1,2,l}$ per-layer \\
Beam $k>1$ & offset cố định $(k_1,k_2)$ từ Bảng~5.2.2.2.1a-3 & mỗi tầng chọn độc lập \\
Không gian ứng viên (hạng~2) & $16{,}384$ & $\sim 10^6$ (không tính trước được) \\
Lợi thế & Đơn giản, codebook tính trước & Bám sát kênh multi-rank phân tán \\
\bottomrule
\end{tabular}
\end{table}

\subsection{Năm thay đổi trong \texttt{nrPMIReport.m}}

\begin{enumerate}

\item \textbf{Ghi đè số cổng (dòng~$\sim$375)}: Code gốc đọc số cổng từ
\texttt{csirs.NumCSIRSPorts} = 32 (một tài nguyên). Với 128 cổng tổng hợp,
phải lấy từ chiều của ma trận kênh:
\begin{lstlisting}[style=matlab]
if strcmpi(reportConfig.CodebookType,'typeI-SinglePanel-r19')
    numCSIRSPorts = size(H,4);   % 128, not 32
end
\end{lstlisting}

\item \textbf{Flag \texttt{isType1SinglePanel\_r19} và \texttt{isModeB\_r19} (dòng~$\sim$386)}:
\begin{lstlisting}[style=matlab]
isType1SinglePanel_r19 = strcmpi(reportConfig.CodebookType,'typeI-SinglePanel-r19');
if isType1SinglePanel_r19; isType1SinglePanel = true; end
isModeB_r19 = isType1SinglePanel_r19 && (reportConfig.CodebookMode == 2);
\end{lstlisting}
Flag \texttt{isModeB\_r19} phân nhánh luồng thực thi sớm, trước khi code
cũ cố dùng codebook tensor (vốn không tồn tại cho Mode~B).

\item \textbf{Nhánh tính PMI Mode~B (dòng~$\sim$432)}: Gọi trực tiếp hàm
greedy search và bỏ qua toàn bộ luồng tra cứu codebook cũ:
\begin{lstlisting}[style=matlab]
if isModeB_r19
    nVar_scalar = mean(nVar(:), 'omitnan');
    [W, pmiModeB, sinrVec] = getPMIType1SinglePanelR19CodebookModeB(...
        Hcsirs, reportConfig, nLayers, nVar_scalar, numCSIRSPorts);
    PMISet.i1    = [pmiModeB.q1+1, pmiModeB.q2+1, pmiModeB.n_idx+1];
    PMISet.i2    = pmiModeB.c_idx + 1;
    SINRPerREPMI = sinrVec(:).';
end
\end{lstlisting}

\item \textbf{Bỏ qua xử lý subband và gán output cho Mode~B (dòng~$\sim$462--471)}:
Mode~B chỉ hỗ trợ wideband, nên hai khối subband và form-output được bao bởi
điều kiện \texttt{if \textasciitilde isModeB\_r19}.

\item \textbf{Placeholder trong \texttt{getCodebook()} (dòng~$\sim$476/493)}: Sau
\texttt{getCodebook()}, code gốc kiểm tra \texttt{\textasciitilde any(codebook(:))}
và thoát sớm nếu codebook rỗng. Mode~A trả về tensor đầy đủ; Mode~B trả về
\texttt{complex(1)} (scalar $\neq 0$) để vượt qua kiểm tra này mà không làm sai
các luồng khác.

\end{enumerate}

\subsection{Hàm Mode A: \texttt{getPMIType1SinglePanelR19CodebookModeA()}}

Mode~A dùng chiến lược \emph{tính trước codebook}. Toàn bộ bộ tiền mã hoá
ứng viên được tạo ra một lần dưới dạng tensor 6 chiều, sau đó hàm tìm PMI
chọn codeword cho SINR cao nhất bằng exhaustive search.

Với $(N_1,N_2)=(16,4)$, $(O_1,O_2)=(4,4)$, số ứng viên ở hạng~1 là:
$N_{i_2} \times N_{i_{1,1}} \times N_{i_{1,2}} \times N_{i_{1,3}} = 4 \times 64 \times 16 \times 4 = 4{,}096$.
Kích thước tensor ở các hạng cao hơn vẫn khả thi vì ràng buộc beam cố định
giữ không gian tìm kiếm tuyến tính theo số ứng viên offset $(k_1,k_2)$.

\subsection{Hàm Mode B: Thuật toán Greedy Factored Search}

Mode~B cho phép mỗi tầng chọn beam DFT \emph{hoàn toàn độc lập}, điều này
mang lại hiệu năng cao hơn trên kênh multi-rank có singular vectors phân tán,
nhưng khiến không gian ứng viên tăng theo lũy thừa: hạng~2 yêu cầu $\sim$4,3\,GB
tensor --- không khả thi để tính trước như Mode~A.

Giải pháp là thuật toán \textbf{Greedy Factored Search} chia bài toán thành
hai lớp tối ưu hoá:

\paragraph{Lớp 1 --- Chọn offset lưới DFT.}
Mode~B sử dụng \emph{offset phân số} $(q_1, q_2)$ để xác định lưới con DFT
cho cả slot. Có $O_1 \times O_2 = 16$ cặp offset. Thuật toán lặp qua cả 16 cặp.

\paragraph{Lớp 2 --- Chọn beam và co-phase.}
Với mỗi offset đã chọn, beam tốt nhất cho từng tầng được xác định:
\begin{itemize}[noitemsep]
  \item Hạng~1: duyệt vét cạn cả $N_1 N_2 = 64$ beam trong lưới con
        (đảm bảo tối ưu, bằng Mode~A).
  \item Hạng~$v > 1$: dùng SVD của $\mathbf{H}_{wb}$ để chọn greedy 1 beam
        tốt nhất cho mỗi tầng. Điểm số phải tính trên \emph{cả hai phân cực}:
        \begin{equation}
          \text{score}(n,l)
          = \underbrace{|\mathbf{V}_q^H\mathbf{V}_\text{top}|_{n,l}^2}_{\text{phân cực trên}}
          + \underbrace{|\mathbf{V}_q^H\mathbf{V}_\text{bot}|_{n,l}^2}_{\text{phân cực dưới}}
        \end{equation}
        Bỏ qua $\mathbf{V}_\text{bot}$ gây mất $\approx50\%$ thông tin trên kênh Rayleigh.
\end{itemize}
Sau khi có bộ beam, duyệt vét cạn $4^v$ tổ hợp co-phase và chọn $W$ tối đa hoá
dung lượng Shannon.

\paragraph{Độ phức tạp.}
Hạng~4: $16\times(64+256)=5{,}120$ đánh giá thay vì
$\approx 6{,}9\times10^{10}$ (vét cạn toàn bộ).

%% ═══════════════════════════════════════════════════════════════════════════
\FloatBarrier
\section{Mở rộng Thư viện: \texttt{nrCQISelect} cho Rel-19}
%% ═══════════════════════════════════════════════════════════════════════════

Bản sao thư viện MATLAB 5G Toolbox cục bộ trong dự án
(\texttt{5g/netmod/+nr5g/+internal/nrCQISelect.m}) chưa hỗ trợ kiểu codebook
\texttt{'typeI-SinglePanel-r19'}. Năm thay đổi có mục tiêu đã được thực hiện:

\begin{enumerate}
    \item \textbf{Xác thực đầu vào} (dòng~$\sim$1018): Thêm
          \texttt{'typeI-SinglePanel-r19'} vào danh sách \texttt{validatestring}
          cho \texttt{CodebookType}.

    \item \textbf{Giới hạn số lớp} (dòng~$\sim$1102): Thêm nhánh \texttt{elseif}
          cho Rel-19 thiết lập \texttt{codebookType='Rel-19 Refined Type I'; maxNLayers=4}.

    \item \textbf{Nhánh chọn PMI} (dòng~$\sim$616): Khi
          \texttt{CodebookType='typeI-SinglePanel-r19'}, hàm giờ gọi
          \texttt{nrPMIReport} với \texttt{nrCSIReportConfig} được xây dựng thủ công
          thay vì \texttt{nrDLPMISelect} cũ:

\begin{lstlisting}[style=matlab, caption={Nhánh PMI thêm vào nrCQISelect cho Rel-19}]
if strcmpi(reportConfig.CodebookType,'typeI-SinglePanel-r19')
    r19cfg = nrCSIReportConfig;
    r19cfg.NSizeBWP          = carrier.NSizeGrid;
    r19cfg.CodebookType      = 'typeI-SinglePanel-r19';
    r19cfg.PanelDimensions   = [1, 16, 4];
    r19cfg.PMIFormatIndicator = 'wideband';
    r19cfg.CodebookMode      = reportConfig.CodebookMode;
    [PMISet, PMIInfo] = nr5g.internal.nrPMIReport(carrier, csirs, ...
                            r19cfg, nLayers, H, nVar);
else
    [PMISet, PMIInfo] = nr5g.internal.nrDLPMISelect(...);
end
\end{lstlisting}

    \item \textbf{Ghi đè số cổng} (dòng~$\sim$1132): Với Rel-19, số cổng CSI-RS
          dùng để xác thực chiều H được ghi đè thành \texttt{size(H,4)}
          (số anten TX thực tế = 128) thay vì \texttt{csirs\{1\}.NumCSIRSPorts} (= 32).

    \item \textbf{Mở rộng chiều SINR} (dòng~$\sim$649): \texttt{nrPMIReport} trả về
          \texttt{SINRPerREPMI} dạng $[1 \times N_L]$ (scalar wideband), nhưng
          \texttt{getSINRperRB} phía sau yêu cầu $[N_{RE} \times N_L]$. Cách sửa
          là broadcast:
\begin{lstlisting}[style=matlab]
if strcmpi(reportConfig.CodebookType,'typeI-SinglePanel-r19') ...
        && size(sinrPerREPMI,1)==1
    sinrPerREPMI = repmat(sinrPerREPMI, length(csirsIndSubs_k), 1);
    PMIInfo.SINRPerREPMI = sinrPerREPMI;
end
\end{lstlisting}
\end{enumerate}

%% ═══════════════════════════════════════════════════════════════════════════
\FloatBarrier
\section{Mở rộng Thư viện: \texttt{nrRISelect} cho Rel-19}
%% ═══════════════════════════════════════════════════════════════════════════

Hàm \texttt{nr5g.internal.nrRISelect} trong thư viện cục bộ chọn RI bằng cách
lặp qua các rank ứng viên và gọi \texttt{nrDLPMISelect} --- một hàm chỉ hỗ trợ
Rel-15/16. Do đó, \textbf{7 thay đổi có mục tiêu} đã được thực hiện để thêm
nhánh Rel-19 mà không ảnh hưởng đến luồng hiện có.

\subsection{Tại sao nâng cấp tại \texttt{nrRISelect}, không phải \texttt{nrDLPMISelect}}

\texttt{nrRISelect} gọi một trong hai hàm nội bộ:
\begin{itemize}[noitemsep]
  \item \texttt{riSelectPMI}: tiêu chí MaxSINR --- gọi \texttt{nrDLPMISelect}, sau đó
        index kết quả qua \texttt{PMIInfo.SINRPerSubband(idx,:,\,PMI.i2(idx),\,PMI.i1(1),\ldots)}.
  \item \texttt{riSelectCQI}: tiêu chí MaxSE --- gọi \texttt{nrCQISelect}.
\end{itemize}

Câu hỏi tự nhiên đặt ra: \textbf{tại sao không mở rộng \texttt{nrDLPMISelect} theo
triết lý nâng cấp toolbox?}

Refined Codebook Rel-19 \emph{có} hỗ trợ subband (TS~38.214 \S5.2.2.2.1a quy định
$i_2$ co-phase có thể báo cáo theo từng subband), nên lập luận ``Rel-19 không có
subband'' là không chính xác. Vấn đề thực sự nằm ở \textbf{cấu trúc output
không tương thích}: sau khi \texttt{nrDLPMISelect} trả về, \texttt{riSelectPMI}
index thẳng vào tensor SINR bằng các chỉ số codebook Rel-15/16:
\begin{lstlisting}[style=matlab]
% riSelectPMI, dong ~439
subbandSINRs(idx,:) = PMIInfo.SINRPerSubband(idx,:, ...
    PMI.i2(idx), PMI.i1(1), PMI.i1(2), PMI.i1(3)) * rankIdx;
\end{lstlisting}
Rel-19 dùng cấu trúc chỉ số $[i_2,\,i_{1,1},\,i_{1,2},\,i_{1,3}]$ (Mode~A)
hoặc $[q_1,\,q_2,\,n^{(1)},\ldots,n^{(v)},\,c^{(1)},\ldots,c^{(v)}]$ (Mode~B)
--- không khớp với layout tensor Rel-15/16. Nếu mở rộng \texttt{nrDLPMISelect},
nó buộc phải trả về các trường \texttt{i1}/\texttt{i2}/\texttt{SINRPerSubband}
giả tạo, trong khi bên trong thực ra đang gọi \texttt{nrPMIReport} với tiêu chí
capacity hoàn toàn khác. Đó là overhead vô nghĩa.

Cách nâng cấp đúng tầng là tại \textbf{\texttt{nrRISelect}} --- hàm quyết định RI.
\texttt{nrDLPMISelect} và \texttt{nrPMIReport} là hai implementation song song,
mỗi cái phục vụ một thế hệ codebook:

\begin{center}
\begin{tabular}{ll}
\toprule
\textbf{Hàm} & \textbf{Vai trò} \\
\midrule
\texttt{nrDLPMISelect} & Sinh codebook + chọn PMI cho Rel-15/16 \\
\texttt{nrPMIReport}   & Chọn PMI cho Rel-19 (128 cổng, Mode A/B) \\
\texttt{nrRISelect}    & Điều phối RI --- gọi đúng hàm tuỳ theo codebook type \\
\bottomrule
\end{tabular}
\end{center}

Nhánh \texttt{riSelectR19} thêm vào \texttt{nrRISelect} chính là điểm nâng cấp
đúng tầng: khi phát hiện Rel-19, route thẳng đến \texttt{nrPMIReport},
bỏ qua \texttt{nrDLPMISelect} (giữ nguyên, không cần chạm vào).
Luồng Rel-15/16 hoàn toàn không bị ảnh hưởng.

\subsection{Bảy thay đổi trong \texttt{nrRISelect.m}}

\begin{enumerate}

\item \textbf{Xác thực \texttt{CodebookType} (dòng~$\sim$586)}: Thêm token mới
vào \texttt{validatestring}:
\begin{lstlisting}[style=matlab]
reportConfig.CodebookType = validatestring(reportConfig.CodebookType,...
    {'Type1SinglePanel','Type1MultiPanel','Type2','eType2',...
     'typeI-SinglePanel-r19'}, fcnName, 'CodebookType field');
\end{lstlisting}

\item \textbf{Giới hạn rank tối đa (dòng~$\sim$368)}: Rel-19 hỗ trợ tối đa
hạng~4 (theo TS~38.214 \S5.2.2.2.1a):
\begin{lstlisting}[style=matlab]
if strcmpi(reportConfig.CodebookType,'typeI-SinglePanel-r19')
    maxRank = min(nRxAnts, 4);
elseif strcmpi(reportConfig.CodebookType,'Type1SinglePanel')
    maxRank = min([nRxAnts Pcsirs 8]);
...
\end{lstlisting}

\item \textbf{Chuỗi mô tả \texttt{codebookType} (dòng~$\sim$657)}: Dùng trong
thông báo lỗi \texttt{RIRestriction}:
\begin{lstlisting}[style=matlab]
if strcmpi(reportConfig.CodebookType,'typeI-SinglePanel-r19')
    maxRank = 4;
    codebookType = 'Rel-19 refined type I single-panel';
...
\end{lstlisting}

\item \textbf{Ghi đè số cổng cho xác thực H (dòng~$\sim$708)}: Tương tự
\texttt{nrCQISelect}, \texttt{csirs.NumCSIRSPorts} = 32 nhưng H có 128 cột:
\begin{lstlisting}[style=matlab]
if strcmpi(reportConfig.CodebookType,'typeI-SinglePanel-r19')
    NumCSIRSPorts = size(H,4);
else
    NumCSIRSPorts = csirs.NumCSIRSPorts(1);
end
validateattributes(H,{class(H)},{'size',[K L NaN NumCSIRSPorts]},fcnName,'H');
\end{lstlisting}

\item \textbf{Khởi tạo output \texttt{initOutputs} (dòng~$\sim$510)}: Placeholder
PMISet phù hợp với định dạng \texttt{nrPMIReport} (trường \texttt{W} thay vì
\texttt{i1}/\texttt{i2}):
\begin{lstlisting}[style=matlab]
if strcmpi(reportConfig.CodebookType,'typeI-SinglePanel-r19')
    PMISet = struct('W', []);
...
\end{lstlisting}

\item \textbf{Phân nhánh chính (dòng~$\sim$398)}: Gọi hàm mới thay vì
\texttt{riSelectPMI}/\texttt{riSelectCQI}:
\begin{lstlisting}[style=matlab]
if ~isempty(validRanks) && ~isempty(csirsInd)
    if strcmpi(reportConfig.CodebookType,'typeI-SinglePanel-r19')
        [RI,PMISet,PMIInfo] = riSelectR19(carrier,csirs,reportConfig,...
                                          H,nVar,validRanks);
    elseif strcmpi(alg,'MaxSINR')
        [RI,PMISet,PMIInfo] = riSelectPMI(...);
    else
        [RI,PMISet,PMIInfo] = riSelectCQI(...);
    end
end
\end{lstlisting}

\item \textbf{Hàm mới \texttt{riSelectR19()} (cuối file)}: Lặp qua
\texttt{validRanks}, gọi \texttt{nrPMIReport} cho từng rank, chọn rank tối đa
hoá dung lượng Shannon:
\begin{lstlisting}[style=matlab]
function [RI,PMISet,PMIInfo] = riSelectR19(carrier,csirs,reportConfig,H,nVar,validRanks)
    H_wb = reshape(mean(mean(H,1),2), size(H,3), size(H,4));
    nRx  = size(H_wb,1);
    r19cfg = nrCSIReportConfig;
    r19cfg.NSizeBWP        = reportConfig.NSizeBWP;
    r19cfg.NStartBWP       = reportConfig.NStartBWP;
    r19cfg.CodebookType    = 'typeI-SinglePanel-r19';
    r19cfg.PanelDimensions = reportConfig.PanelDimensions;
    r19cfg.PMIFormatIndicator = 'wideband';
    r19cfg.CodebookMode    = reportConfig.CodebookMode;
    bestCap = -Inf;
    for rank = validRanks
        try
            [pmiSet_try,pmiInfo_try] = nr5g.internal.nrPMIReport(...
                carrier,csirs,r19cfg,rank,H,nVar);
            W   = pmiInfo_try.W;
            HW  = H_wb * W;
            cap = real(log2(det(eye(nRx) + (1/nVar)*(HW*HW'))));
            if cap > bestCap
                bestCap = cap;  RI = rank;
                PMISet = pmiSet_try;  PMIInfo = pmiInfo_try;
            end
        catch; end
    end
end
\end{lstlisting}

\end{enumerate}

\subsection{Tích hợp vào \texttt{csirs\_feedback.m}}

Sau khi mở rộng \texttt{nrRISelect}, các vòng lặp thủ công trong
\texttt{csirs\_feedback.m} được thay bằng:

\begin{lstlisting}[style=matlab]
% Approach C (Mode A)
riCfg_C.NSizeBWP       = carrier.NSizeGrid;
riCfg_C.NStartBWP      = 0;
riCfg_C.CodebookType   = 'typeI-SinglePanel-r19';
riCfg_C.PanelDimensions = [1, 16, 4];
riCfg_C.CodebookMode   = 1;
[ri_C, ~, ~] = nr5g.internal.nrRISelect(carrier, csirs{1}, riCfg_C, H_4d, nVar_wb);

% Approach D (Mode B): riCfg_D identique avec CodebookMode = 2
[ri_D, ~, ~] = nr5g.internal.nrRISelect(carrier, csirs{1}, riCfg_D, H_4d, nVar_wb);
\end{lstlisting}

Toàn bộ quá trình chọn RI giờ đi qua đúng hàm toolbox
\texttt{nrRISelect} $\to$ \texttt{riSelectR19} $\to$ \texttt{nrPMIReport},
đảm bảo nhất quán với triết lý \emph{tối đa hoá 5G Toolbox, mở rộng có mục tiêu}.

%% ═══════════════════════════════════════════════════════════════════════════
\FloatBarrier
\section{Lịch sử Debug}
%% ═══════════════════════════════════════════════════════════════════════════

Các lỗi sau đây đã phát sinh và được giải quyết trong quá trình phát triển v2.0:

\begin{table}[h!]
\centering
\caption{Các lỗi gặp phải và cách khắc phục}
\label{tab:errors}
\begin{tabular}{p{4cm}p{5cm}p{5cm}}
\toprule
\textbf{Lỗi} & \textbf{Nguyên nhân gốc} & \textbf{Cách khắc phục} \\
\midrule
\texttt{Expected H to be of size ...x32} &
\texttt{parseInputs} đặt \texttt{nTxAnts=32} từ tài nguyên CSI-RS (không phải 128 đầy đủ) &
Ghi đè \texttt{numCSIRSPorts = size(H,4)} cho Rel-19 trong \texttt{validateInputs} \\
\addlinespace
\texttt{Index exceeds array bounds in getSINRperRB} &
\texttt{nrPMIReport} trả về \texttt{SINRPerREPMI [1×nL]}, phía sau kỳ vọng \texttt{[nRE×nL]} &
Broadcast SINR wideband sang per-RE qua \texttt{repmat} sau khi chọn PMI \\
\addlinespace
\texttt{Unrecognized function nrCSIReferenceResource} &
Hàm này tồn tại trong toolbox MATLAB chính thức nhưng không có trong bản sao 5g/ cục bộ &
Thay bằng \texttt{nrPDSCHConfig} xây dựng thủ công (PRBSet + NumLayers) \\
\addlinespace
\texttt{Unrecognized field name sinrPerRE\_B} &
\texttt{csirs\_feedback.m} không lưu \texttt{sinr\_B} vào struct \texttt{fb} &
Thêm \texttt{fb.sinr\_B = sinr\_B;} và cập nhật \texttt{csirs\_computeMetrics} dùng \texttt{repmat} \\
\addlinespace
\texttt{Unrecognized property 'MCSTable' for nrPDSCHConfig} &
\texttt{nrPDSCHConfig} dùng thuộc tính \texttt{Modulation}, không phải \texttt{MCSTable} &
Xóa \texttt{pdsch.MCSTable = 'qam64'} khỏi hàm trợ giúp \texttt{makePDSCH} \\
\addlinespace
\texttt{Mode B trả về NaN (không tính PMI)} &
\texttt{getCodebook()} trả về mảng rỗng, kích hoạt kiểm tra \texttt{\textasciitilde any(codebook(:))} và thoát sớm &
Trả về \texttt{complex(1)} làm placeholder --- scalar phức khác~0 vượt qua NaN-check mà không làm sai luồng khác \\
\addlinespace
\texttt{Mode B TEST 5 fail (Mode B thua Mode A 85\%)} &
Điểm số beam chỉ dùng phân cực trên ($V_\text{top}$), bỏ sót 50\% năng lượng của kênh Rayleigh &
Cộng thêm đóng góp phân cực dưới ($V_\text{top} + V_\text{bot}$) vào hàm tính điểm greedy \\
\addlinespace
\texttt{Mode B TEST 5 vẫn fail sau fix scoring} &
Hạng~1 dùng greedy top-1 (64 lần ít hơn Mode A), không đảm bảo tối ưu &
Thay \texttt{n\_cand = top1} thành \texttt{n\_cand = 0:N1N2-1} (vét cạn) khi \texttt{nLayers == 1}; TEST 5 đạt 100\% \\
\addlinespace
\texttt{function not allowed in if/else block} &
Local function không thể đặt bên trong khối điều kiện MATLAB script &
Inline toàn bộ vòng lặp RI trực tiếp vào thân script; sau này chuyển thành hàm riêng trong \texttt{nrRISelect} \\
\bottomrule
\end{tabular}
\end{table}

%% ═══════════════════════════════════════════════════════════════════════════
\FloatBarrier
\section{Kết quả Mô phỏng}
%% ═══════════════════════════════════════════════════════════════════════════

\subsection{Kết quả thô (20 realization mỗi SNR)}

\begin{table}[h!]
\centering
\caption{Kết quả Quét SNR: RI trung bình, MCS trung bình, Thông lượng (Mbps), \% so với cận trên SVD}
\label{tab:results}
\begin{tabular}{llccc}
\toprule
\textbf{SNR} & \textbf{Phương pháp} & \textbf{RI TB} & \textbf{TP trung bình (Mbps)} & \textbf{\% so với B} \\
\midrule
\multirow{3}{*}{$-5$\,dB (xa)}
  & B: SVD (tham chiếu) & --- & 216{,}2 & tham chiếu \\
  & C: Mode A           & --- & 59{,}3  & 27{,}4\,\% \\
  & D: Mode B           & --- & 72{,}1  & 33{,}4\,\% \\
\addlinespace
\multirow{3}{*}{$+10$\,dB (trung bình)}
  & B: SVD (tham chiếu) & --- & 332{,}7 & tham chiếu \\
  & C: Mode A           & --- & 181{,}5 & 54{,}5\,\% \\
  & D: Mode B           & --- & 198{,}1 & 59{,}5\,\% \\
\addlinespace
\multirow{3}{*}{$+25$\,dB (gần)}
  & B: SVD (tham chiếu) & --- & 332{,}7 & tham chiếu \\
  & C: Mode A           & --- & 216{,}2 & 65{,}0\,\% \\
  & D: Mode B           & --- & 233{,}8 & 70{,}3\,\% \\
\bottomrule
\end{tabular}
\end{table}

\subsection{Nhận xét}

\begin{enumerate}
    \item \textbf{Mode~B luôn nhỉnh hơn Mode~A.}
          Chênh lệch khoảng 5--6\,\% so với SVD ở mọi mức SNR. Không nhiều nhưng
          nhất quán, phù hợp với kỳ vọng vì Mode~B dùng cơ sở trực giao nên bám
          sát hơn với cấu trúc mảng 2D $(N_1=16, N_2=4)$.

    \item \textbf{SNR thấp thì cả hai đều xa SVD.}
          Ở $-5$\,dB, Mode~A chỉ đạt 27\,\%, Mode~B 33\,\%. Khi SNR tăng lên
          $+25$\,dB, khoảng cách thu hẹp rõ rệt (65\,\% và 70\,\%). Phần còn lại
          là do lượng hóa codebook -- không thể tránh được.

    \item \textbf{Thông lượng SVD trần ở $+10$\,dB và $+25$\,dB giống nhau.}
          Cả hai điểm đều cho $\approx 332{,}7$\,Mbps. Lý do là MCS đã chạm trần
          (QAM256, $R \approx 0{,}93$) từ $+10$\,dB với cấu hình 10\,MHz này rồi.

    \item \textbf{RI giảm khi delay spread và Doppler cao -- không phải lỗi codebook.}
          Với CDL-B ở 100\,ns / 5\,Hz Doppler thì kênh rất ổn, dễ ước lượng, RI
          thường đạt 4. Khi tăng lên 450\,ns / 136\,Hz, chất lượng ước lượng kênh
          giảm, $\mathbf{H}_{wb}$ bị trải ra trên miền tần số và thời gian, rank
          hiệu dụng thấp đi. Đây là vấn đề của kênh và CE, không phải codebook bị
          hạn chế.
\end{enumerate}

%% ═══════════════════════════════════════════════════════════════════════════
\FloatBarrier
\section{Tóm tắt Thay đổi: v1.0 đến v2.0}
%% ═══════════════════════════════════════════════════════════════════════════

\begin{table}[h!]
\centering
\caption{So sánh: v1.0 và v2.0}
\label{tab:comparison}
\begin{tabular}{lll}
\toprule
\textbf{Khía cạnh} & \textbf{v1.0} & \textbf{v2.0} \\
\midrule
Cấu trúc mã          & Script đơn khối 12 mục        & 7 file hàm + 1 file chính \\
Quét SNR             & Một điểm (20\,dB)             & 3 điểm ($-5$, $+10$, $+25$\,dB) \\
Realization kênh     & 1 (tất định)                  & 20 mỗi SNR (Monte Carlo qua \texttt{release}) \\
Các phương pháp      & A, B, C, D                    & B, C, D (bỏ A vì là legacy Rel-15/16) \\
Tính CQI             & Bảng ngưỡng SINR              & L2SM/MIESM qua \texttt{nrCQISelect} \\
Tính MCS             & Dùng chỉ số CQI làm MCS       & \texttt{computeMCS} (L2SM, TS~38.214) \\
Chọn RI (C/D)        & Vòng lặp \texttt{for ri\_try} thủ công & \texttt{nr5g.internal.nrRISelect} (7 thay đổi) \\
\texttt{nrCQISelect} Rel-19 & Không hỗ trợ & Mở rộng với 5 thay đổi có mục tiêu \\
\texttt{nrPMIReport} Rel-19 & Mode A (32p, không có Mode B) & Mode A 128p + Mode B Greedy Factored Search \\
Đầu ra               & In một điểm                   & Bảng ASCII với \% so với SVD \\
\bottomrule
\end{tabular}
\end{table}

%% ═══════════════════════════════════════════════════════════════════════════
\section{Kết luận}
%% ═══════════════════════════════════════════════════════════════════════════

Từ v1.0 đến v2.0, pipeline mô phỏng 128T128R đã đi từ một script chạy thử một lần
thành một công cụ đánh giá có thể chạy lại được, với kết quả thống kê tin cậy.
Tóm lại những gì đã làm được:

\begin{itemize}
    \item Tách code thành 7 file hàm rõ ràng, dễ bảo trì và tái sử dụng.
    \item Mở rộng \texttt{nrPMIReport}: cài đặt Mode~A (128 cổng, tensor codebook DFT 2D)
          và Mode~B (Greedy Factored Search, $\leq$5{,}120 đánh giá thay vì $\sim10^{10}$),
          xác nhận bằng 5 bài test (tất cả PASS).
    \item Mở rộng \texttt{nrCQISelect} (5 thay đổi) và \texttt{nrRISelect} (7 thay đổi)
          để hỗ trợ \texttt{'typeI-SinglePanel-r19'}, đảm bảo toàn bộ
          pipeline CSI đi qua đúng hàm toolbox.
    \item Chuyển sang tính MCS bằng \texttt{computeMCS} (L2SM/MIESM, TS~38.214
          Bảng~5.1.3.1-1), không còn dùng chỉ số CQI như proxy cho MCS.
    \item Kết quả số (CDL-B, 20 realization/SNR): Mode~B đạt 33--70\,\% so với
          cận trên SVD; Mode~A đạt 27--65\,\%; chênh nhau $\approx 5$--6\,\% nhất quán
          ở mọi SNR.
\end{itemize}

\appendix

%% ═══════════════════════════════════════════════════════════════════════════
\section{Danh sách File}
%% ═══════════════════════════════════════════════════════════════════════════

\begin{tabular}{ll}
\toprule
\textbf{File} & \textbf{Nhánh / Vị trí} \\
\midrule
\texttt{CSIRS\_128T128R\_v1\_0.m}              & \texttt{main} \\
\texttt{CSIRS\_128T128R\_v2\_0.m}              & \texttt{SNRSweep} \\
\texttt{csirs\_buildGrid.m}                    & \texttt{SNRSweep} \\
\texttt{csirs\_channelSetup.m}                 & \texttt{SNRSweep} \\
\texttt{csirs\_txRx.m}                         & \texttt{SNRSweep} \\
\texttt{csirs\_channelEstimate.m}              & \texttt{SNRSweep} \\
\texttt{csirs\_feedback.m}                     & \texttt{SNRSweep} \\
\texttt{csirs\_computeMetrics.m}               & \texttt{SNRSweep} \\
\texttt{csirs\_printSweepTable.m}              & \texttt{SNRSweep} \\
\texttt{5g/netmod/+nr5g/+internal/nrCQISelect.m} & \texttt{SNRSweep} (đã sửa đổi, 5 thay đổi Rel-19) \\
\texttt{5g/netmod/+nr5g/+internal/nrRISelect.m} & \texttt{SNRSweep} (đã sửa đổi, 7 thay đổi Rel-19) \\
\texttt{5g/netmod/+nr5g/+internal/nrPMIReport.m} & \texttt{SNRSweep} (đã sửa đổi, 5 thay đổi + 2 hàm mới) \\
\texttt{test\_ModeB\_verification.m} & \texttt{SNRSweep} (5 bài test, tất cả PASS) \\
\bottomrule
\end{tabular}

%% ═══════════════════════════════════════════════════════════════════════════
\section{Tài liệu tham khảo 3GPP}
%% ═══════════════════════════════════════════════════════════════════════════

\begin{itemize}[noitemsep]
    \item TS~38.211~§7.4.1.5 --- Tạo tín hiệu NZP-CSI-RS
    \item TS~38.214~§5.2.2.2.1 --- Codebook Type~I Single-Panel (Rel-15/16)
    \item TS~38.214~§5.2.2.2.1a --- Codebook Refined Type~I Single-Panel (Rel-19)
    \item TS~38.214~§5.2.2.1 --- Báo cáo CQI và ràng buộc BLER
    \item TS~38.214 Bảng~5.1.3.1-1 --- Bảng chỉ số MCS (64QAM)
    \item TS~38.331 --- Cấu hình RRC NZP-CSI-RS-ResourceSet
\end{itemize}

\end{document}
